% !TEX root = ../main.tex
\chapter{同型：安全な相互変換}
\label{chap:ch12-isomorphism}

\begin{chapterabstract}
本章では、同型を「情報を失わない相互変換」として導入する。ソフトウェアでは、スキーマ移行、DTO変換、encoder / decoder、旧APIと新APIの対応などで、形は変わるが意味は変えたくない場面が多い。ただし、$A \to B$ と $B \to A$ という二つの関数があるだけでは十分ではない。変換して戻すと元に戻ること、逆向きにも戻ること、つまり二つの round-trip law が必要である。本章では、この法則を小さなLean構造として定義し、\texttt{UserV1} から \texttt{UserV2} への lossless migration を題材に確認する。
\end{chapterabstract}

\begin{learninggoals}
\item 同型を、単なる「似ている」ではなく「情報を失わない相互変換」として説明できる。
\item $A \to B$ と $B \to A$ の存在だけでは同型にならない理由を説明できる。
\item $\mathit{from} \circ \mathit{to} = \mathit{id}$ と $\mathit{to} \circ \mathit{from} = \mathit{id}$ の意味を読める。
\item Lean 4で、\texttt{toFun}、\texttt{invFun}、二つの証明を持つ小さな同型構造を定義できる。
\item スキーマ移行を完全同型、片方向互換、情報損失ありに分けて考えられる。
\end{learninggoals}

\section{この章を読む理由}
\label{sec:ch12-why}

ソフトウェア開発では、形を変えたいが意味は変えたくない場面が多い。たとえば、次のような変更である。

\begin{itemize}
\item \texttt{UserV1} の平坦なフィールドを、\texttt{UserV2} では \texttt{contact} と \texttt{profile} に分ける。
\item 外部APIから届くJSONを内部DTOに変換し、必要なら再びJSONへ戻す。
\item DBスキーマのカラム名を変えるが、保持している情報は変えない。
\item 古いレスポンス形式を新しいレスポンス形式へadapterで変換する。
\end{itemize}

このようなとき、私たちはよく「互換性がある」「戻せる」「意味は同じ」と言う。しかし、その言葉だけでは曖昧である。何をどの向きに戻せるのか。どの範囲のデータなら戻せるのか。情報を落としていないのか。そこを明確にするための言葉が、同型である。

図\ref{fig:ch12-overview}では、正方向と逆方向の変換に加え、両方の合成がそれぞれの恒等変換になるという二つの証明を分けて示す。

\FigurePlaceholder{fig-ch12-overview.png}{0.92}{同型のデータと法則}{fig:ch12-overview}

\section{同型は「情報を失わない変換」}
\label{sec:ch12-information-preserving}

第11章では、型を対象、関数を射として見た。そこでは、関数は一方向の変換だった。本章では、二つの型の間に、互いに戻し合う二つの関数がある場合を考える。

\begin{definitionbox}
型 $A$ と型 $B$ が同型であるとは、次のものがあるということである。

\begin{itemize}
\item $A$ から $B$ への変換 $\mathit{to} : A \to B$
\item $B$ から $A$ への変換 $\mathit{from} : B \to A$
\item $A$ から $B$ へ行って $A$ に戻ると元に戻るという証明
\item $B$ から $A$ へ行って $B$ に戻ると元に戻るという証明
\end{itemize}

ここで重要なのは、最後の二つが「実装上の期待」ではなく、法則として要求されることである。
\end{definitionbox}

同型は、「同じ型である」という意味ではない。\texttt{UserV1} と \texttt{UserV2} は別の型でよい。フィールド名やネスト構造が違っても、片方の情報をもう片方で完全に表せればよい。なお「同じ情報」という言い方は型と関数の圏での直感であり、一般の圏での正式な条件は、二つの射の合成が両側の恒等射になることである。

たとえば、次の二つは見た目が違う。

\begin{itemize}
\item \texttt{UserV1}: \texttt{id}, \texttt{email}, \texttt{displayName} を直接持つ。
\item \texttt{UserV2}: \texttt{userId}, \texttt{contact.email}, \texttt{profile.displayName} に分けて持つ。
\end{itemize}

この変更は、単にフィールドを移動しているだけなら情報を失わない。旧形式から新形式へ変換し、新形式から旧形式へ戻すことができる。逆向きも同じである。このとき、二つの型は同型として扱える。

\section{encoder / decoder と round-trip}
\label{sec:ch12-round-trip}

同型の感覚は、encoder / decoder の例でつかみやすい。たとえば、アプリ内の値をJSON文字列へ変換する関数を \texttt{encode}、JSON文字列から値へ戻す関数を \texttt{decode} とする。

以下の二式は、\texttt{decode : B -> A} が全域関数である理想化、または (B) を復号可能な表現だけに制限した型として読む。実際の parser が \texttt{Option A} や \texttt{Except E A} を返すなら、そのままでは型の同型ではなく、成功する入力の範囲と失敗時仕様を別に定める必要がある。

数式として書くと、次の二つが必要になる。

\[
  \mathit{decode}(\mathit{encode}(a)) = a
\]

\[
  \mathit{encode}(\mathit{decode}(b)) = b
\]

一つ目は、値側の復元である。二つ目は、表現側の復元である。実務では、一つ目だけをテストして「round-tripできている」と言ってしまうことがある。しかし、完全な同型として見るなら、二つ目も必要になる。

図\ref{fig:ch12-round-trip}では、値側と表現側のどちらから出発しても、往復後に出発点へ戻る二つの経路を区別する。

\FigurePlaceholder{fig-ch12-round-trip.png}{0.90}{二方向の round-trip}{fig:ch12-round-trip}

\begin{warningbox}
JSONやCSVのような外部表現では、表現側に「余分な形」が存在することがある。たとえば、同じ値を表すJSON文字列が空白やキー順序の違いで複数存在する場合、文字列として完全に同じものへ戻るとは限らない。この場合は、完全同型ではなく、正規化後の一致や well-formed な表現に限定した同型として扱う。
\end{warningbox}

\section{\texorpdfstring{$A \to B$ と $B \to A$}{A to B と B to A} だけでは同型ではない}
\label{sec:ch12-bidirectional-is-not-enough}

双方向の関数があることと、同型であることは違う。これは実務で非常に重要である。

極端な例を考える。任意のユーザーを \texttt{Unit} に変換する関数は作れる。\texttt{Unit} からデフォルトユーザーを作る関数も作れる。つまり、\texttt{UserV1} から \texttt{Unit} への関数と、\texttt{Unit} から \texttt{UserV1} への関数は存在する。しかし、これは明らかに安全な相互変換ではない。\texttt{Unit} には、元のメールアドレスや表示名を保持する場所がないからである。

\begin{listing}[htbp]
\begin{minted}{lean4}
structure UserV1 where
  id : Nat
  email : String
  displayName : String
deriving Repr, DecidableEq

def eraseToUnit (_ : UserV1) : Unit := ()

def defaultUser (_ : Unit) : UserV1 :=
  { id := 0, email := "", displayName := "" }

def sampleUser : UserV1 :=
  { id := 1, email := "alice@example.com", displayName := "Alice" }

example : defaultUser (eraseToUnit sampleUser) ≠ sampleUser := by
  decide
\end{minted}
\caption{\texttt{eraseToUnit} と \texttt{defaultUser}}
\label{lst:ch12_bidirectional_not_iso}
\end{listing}

このコードでは、\texttt{eraseToUnit} と \texttt{defaultUser} という二つの向きの関数がある。しかし、\texttt{sampleUser} を \texttt{Unit} に変換して戻すとデフォルトユーザーになり、元のメールアドレスも表示名も戻らない。\texttt{a ≠ b} は二値が等しくないという命題で、\texttt{decide} はこの具体的な不等式を計算で判定する。

図\ref{fig:ch12-bidirectional-not-iso}では、複数のユーザーが同じ \texttt{Unit} に潰れるため、逆向きの関数を一つ選んでも全ユーザーを復元できないことを示す。

\FigurePlaceholder{fig-ch12-bidirectional-not-iso.png}{0.90}{情報を失う双方向変換}{fig:ch12-bidirectional-not-iso}

\section{必要な法則}
\label{sec:ch12-laws}

同型に必要な法則は二つである。型 $A$ から型 $B$ への変換を $\mathit{to}$、型 $B$ から型 $A$ への変換を $\mathit{from}$ とする。

\begin{definitionbox}
同型の二つの法則は次の通りである。

\[
  \mathit{from} \circ \mathit{to} = \mathit{id}_A
\]

\[
  \mathit{to} \circ \mathit{from} = \mathit{id}_B
\]

一つ目は、$A$ の値を $B$ へ変換してから $A$ へ戻すと、元の $A$ の値になるという意味である。二つ目は、$B$ の値を $A$ へ変換してから $B$ へ戻すと、元の $B$ の値になるという意味である。
\end{definitionbox}

ここで、\texttt{id} は恒等関数である。何もしない関数だと考えればよい。つまり、\texttt{from} と \texttt{to} を続けた結果が「何もしない」のと同じである、という主張になる。

実務上は、二つの法則を次のように読む。

\begin{table}[htbp]
\centering
\begin{tabular}{p{0.28\linewidth}p{0.62\linewidth}}
\toprule
法則 & 実務での意味 \\
\midrule
\texttt{rollback (migrate u) = u} & 旧データを新形式へ移行し、戻したときに旧データが完全に復元される。 \\
\texttt{migrate (rollback v) = v} & 新形式のデータを旧形式へ戻し、再び新形式へ移行しても新データが完全に復元される。 \\
\bottomrule
\end{tabular}
\caption{二つの round-trip law}
\label{tab:ch12-round-trip-laws}
\end{table}

\section{Lean 4 で小さな同型構造を定義する}

圏論における同型は、互いに逆となる二つの射を持つこととして定義する~\cite{leinster-basic-category-theory,riehl-category-theory-in-context}。本章の\texttt{Iso}は型と関数の圏に限定した小さなモデルである。
\label{sec:ch12-lean-iso}

本書では、初期章ではMathlibの高度な圏論APIを使わず、自作の小さな構造で概念をつかむ。Leanでは、同型を「二つの関数」と「二つの証明」を持つ構造体として定義できる。

\begin{listing}[htbp]
\begin{minted}{lean4}
universe u v

structure Iso (A : Type u) (B : Type v) where
  toFun : A → B
  invFun : B → A
  left_inv : ∀ a : A, invFun (toFun a) = a
  right_inv : ∀ b : B, toFun (invFun b) = b
\end{minted}
\caption{\texttt{Iso}}
\label{lst:ch12_iso_structure}
\end{listing}

\texttt{universe u v} は、\texttt{A} と \texttt{B} が属する型 universe の段階名を宣言する。これにより、同じ universe に限らず二つの型の間で使える構造になる。この定義は一般の圏の同型を実装したものではないが、「変換関数のペア」に両側の逆法則が必要であることを型と関数の例で明示している。

この構造を使うと、恒等変換は常に同型であると言える。

\begin{listing}[htbp]
\begin{minted}{lean4}
namespace Iso

def refl (A : Type u) : Iso A A where
  toFun := fun a => a
  invFun := fun a => a
  left_inv := by
    intro a
    rfl
  right_inv := by
    intro a
    rfl

end Iso
\end{minted}
\caption{\texttt{Iso.refl}}
\label{lst:ch12_iso_refl}
\end{listing}

\texttt{rfl} は、両辺が定義を展開すると同じ形になるときに使える。ここでは、値をそのまま返しているだけなので、往復しても同じ値になる。

\section{サンプル：\texttt{UserV1} と \texttt{UserV2} の lossless migration}
\label{sec:ch12-user-migration}

次に、実務に近いミニケースを考える。旧形式 \texttt{UserV1} では、ユーザー情報を平坦なレコードとして持つ。新形式 \texttt{UserV2} では、連絡先情報とプロフィール情報をネストした構造として持つ。

図\ref{fig:ch12-user-migration}では、\texttt{id}、\texttt{email}、\texttt{displayName} が新形式のどのフィールドへ一対一に移るかを示す。

\FigurePlaceholder{fig-ch12-user-migration.png}{0.94}{ユーザーフィールドの対応}{fig:ch12-user-migration}

まず、二つの型を定義する。

\begin{listing}[htbp]
\begin{minted}{lean4}
structure UserV1 where
  id : Nat
  email : String
  displayName : String
deriving Repr, DecidableEq

structure Contact where
  email : String
deriving Repr, DecidableEq

structure Profile where
  displayName : String
deriving Repr, DecidableEq

structure UserV2 where
  userId : Nat
  contact : Contact
  profile : Profile
deriving Repr, DecidableEq
\end{minted}
\caption{\texttt{UserV1} と \texttt{UserV2}}
\label{lst:ch12_user_types}
\end{listing}

次に、移行関数とロールバック関数を定義する。

\begin{listing}[htbp]
\begin{minted}{lean4}
def migrate (u : UserV1) : UserV2 :=
  { userId := u.id
    contact := { email := u.email }
    profile := { displayName := u.displayName } }

def rollback (v : UserV2) : UserV1 :=
  { id := v.userId
    email := v.contact.email
    displayName := v.profile.displayName }
\end{minted}
\caption{\texttt{migrate} と \texttt{rollback}}
\label{lst:ch12_user_functions}
\end{listing}

ここまでは、普通のプログラムと変わらない。重要なのは、この二つが本当に情報を失わないことを証明する部分である。

証明方針は単純である。構造体の値を \texttt{cases} で分解し、定義を展開したときに左右が同じ形になることを \texttt{rfl} で示す。

\begin{listing}[htbp]
\begin{minted}{lean4}
theorem rollback_migrate (u : UserV1) :
    rollback (migrate u) = u := by
  cases u
  rfl

theorem migrate_rollback (v : UserV2) :
    migrate (rollback v) = v := by
  cases v with
  | mk userId contact profile =>
    cases contact with
    | mk email =>
      cases profile with
      | mk displayName =>
        rfl
\end{minted}
\caption{\texttt{rollback\_migrate} と \texttt{migrate\_rollback}}
\label{lst:ch12_user_roundtrip}
\end{listing}

一つ目の定理は、旧形式から新形式へ移行して戻すと、旧形式が完全に復元されることを示している。二つ目の定理は、新形式から旧形式へ戻して再び移行しても、新形式が完全に復元されることを示している。

最後に、この二つの関数と二つの証明を \texttt{Iso} にまとめる。

\begin{listing}[htbp]
\begin{minted}{lean4}
def userIso : Iso UserV1 UserV2 where
  toFun := migrate
  invFun := rollback
  left_inv := rollback_migrate
  right_inv := migrate_rollback
\end{minted}
\caption{\texttt{userIso}}
\label{lst:ch12_user_iso}
\end{listing}

\texttt{userIso} は二つの変換だけでなく \texttt{left\_inv} と \texttt{right\_inv} も保持するため、相互変換の法則を後続の証明から再利用できる。

\section{実務上の注意：完全同型・片方向互換・情報損失を区別する}
\label{sec:ch12-practical-classification}

実務では、すべての変換が完全同型になるわけではない。むしろ、多くの変更は、片方向だけ安全だったり、条件付きで戻せたり、情報を落としていたりする。重要なのは、それらを同じ「互換性」という言葉でまとめないことである。

図\ref{fig:ch12-migration-classification}では、二つの round-trip が成り立つ範囲と、失われる情報の有無によって移行を分類する。

\FigurePlaceholder{fig-ch12-migration-classification.png}{0.94}{移行の分類}{fig:ch12-migration-classification}

\begin{table}[htbp]
\centering
\begin{tabular}{p{0.20\linewidth}p{0.31\linewidth}p{0.36\linewidth}}
\toprule
分類 & 成り立つ性質 & 例 \\
\midrule
完全同型 & 両方向のround-tripが成り立つ。 & フィールド名変更、ネスト構造の導入、情報を保つDTO再編成。 \\
片方向互換 & 片方向のround-tripだけが成り立つ。 & 旧データは新形式で読めるが、新形式のすべてを旧形式へ戻せるとは限らない。 \\
制約付き同型 & well-formed 条件の下でだけ両方向に戻せる。 & 正規化済みJSON、キー順序を固定した表現、制約を満たす入力。 \\
情報損失あり & 元に戻せない情報がある。 & 氏名分割、ログ匿名化、集計、丸め、不要フィールドの削除。 \\
\bottomrule
\end{tabular}
\caption{移行ごとの保証}
\label{tab:ch12-migration-classification}
\end{table}

\begin{warningbox}
情報を落とす変更を、無理に同型として扱ってはいけない。証明できない性質を「たぶん戻せる」と書くよりも、「この向きだけ保証する」「この条件の下で保証する」「正規化後の一致だけ保証する」と明示したほうが安全である。
\end{warningbox}

第27章以降では、この分類をより詳しく扱う。第28章では、\texttt{UserV1} と \texttt{UserV2} のような lossless migration をケーススタディとして扱う。第29章では、完全には戻せない lossy migration を扱い、証明可能な仕様へ弱める方法を見る。

\section{設計レビューでの使い方}
\label{sec:ch12-review}

同型の考え方は、Leanコードを書く場面だけでなく、設計レビューでも役に立つ。変換仕様をレビューするときは、次の問いを立てるとよい。

\begin{itemize}
\item 変更前の型と変更後の型は何か。
\item 旧形式から新形式への変換は何か。
\item 新形式から旧形式への変換は何か。
\item \texttt{rollback (migrate x) = x} は成り立つか。
\item \texttt{migrate (rollback y) = y} は成り立つか。
\item 片方だけ成り立つなら、その意味を明記しているか。
\item 条件付きでしか戻せないなら、well-formed 条件を明記しているか。
\item 情報損失があるなら、失われる情報を明記しているか。
\end{itemize}

この問いは、レビューでの会話を具体化する。「互換性がありますか」ではなく、「どちらのround-tripが成り立ちますか」と聞けるようになる。

\section{本章のまとめ}
\label{sec:ch12-summary}

同型は、二つの型が同じ情報を表していることを示す関係である。

$A \to B$ と $B \to A$ の関数があるだけでは同型ではない。両方向のround-trip lawが必要である。

Leanでは、同型を \texttt{toFun}、\texttt{invFun}、\texttt{left\_inv}、\texttt{right\_inv} を持つ構造体として表せる。

\texttt{UserV1} と \texttt{UserV2} のような情報を保つスキーマ変更は、Lean上で同型として表現できる。

実務では、完全同型、片方向互換、制約付き同型、情報損失ありを区別することが重要である。
